% Bibliographie commune : citer uniquement ces clés avec \cite{...}. Ajouter une entrée ici si besoin, jamais ailleurs.
\begin{thebibliography}{99}\setlength{\itemsep}{0pt}
\bibitem{heusel2017} M.~Heusel et al. GANs trained by a two time-scale update rule converge to a local Nash equilibrium. NeurIPS, 2017.
\bibitem{kynkaanniemi2019} T.~Kynkäänniemi et al. Improved precision and recall metric for assessing generative models. NeurIPS, 2019.
\bibitem{kynkaanniemi2023} T.~Kynkäänniemi et al. The role of ImageNet classes in Fréchet Inception Distance. ICLR, 2023.
\bibitem{stein2023} G.~Stein et al. Exposing flaws of generative model evaluation metrics and their unfair treatment of diffusion models. NeurIPS, 2023.
\bibitem{devries2020} T.~DeVries, M.~Drozdzal, G.~W.~Taylor. Instance selection for GANs. NeurIPS, 2020.
\bibitem{sonderby2017} C.~K.~Sønderby et al. Amortised MAP inference for image super-resolution. ICLR, 2017.
\bibitem{arjovsky2017} M.~Arjovsky, L.~Bottou. Towards principled methods for training generative adversarial networks. ICLR, 2017.
\bibitem{jolicoeur2019} A.~Jolicoeur-Martineau. The relativistic discriminator: a key element missing from standard GAN. ICLR, 2019.
\bibitem{mescheder2018} L.~Mescheder, A.~Geiger, S.~Nowozin. Which training methods for GANs do actually converge? ICML, 2018.
\bibitem{huang2024} Y.~Huang, A.~Gokaslan, V.~Kuleshov, J.~Tompkin. The GAN is dead; long live the GAN! A modern GAN baseline. NeurIPS, 2024.
\bibitem{yazici2019} Y.~Yazıcı et al. The unusual effectiveness of averaging in GAN training. ICLR, 2019.
\bibitem{azadi2019} S.~Azadi et al. Discriminator rejection sampling. ICLR, 2019.
\bibitem{verine2024} A.~Verine et al. Optimal budgeted rejection sampling for generative models. AISTATS, 2024.
\end{thebibliography}
